\documentclass[11pt]{article}
\usepackage[a4paper,margin=25mm]{geometry}
\usepackage{amsmath,amssymb,amsthm,microtype,hyperref}
\newtheorem{theorem}{Theorem}\newcommand{\E}{\mathbb E}\newcommand{\Pp}{\mathbb P}
\title{Combining all conditional counts by one crossing identity}
\author{}\date{28 October 2026}
\begin{document}\maketitle

Separate upper bounds on incoming counts discard cancellation with outgoing
counts. The following identity retains every count of the same iid law.
Its leading loss has the correct sign for all noise levels. The proof is a
finite calculation; no Poisson theorem, tail truncation, or parameter grid
enters it. The unequal binomial weights still require a separate estimate.

\begin{theorem}\label{thm:mesh}
Let $U_0,U_1$ be integrable, centered conditional-noise laws with
\[
 U_0\le a_0,\qquad -a_1\le U_1\le A_1,
 \quad a_0,a_1,A_1\ge0,\quad a_0+a_1<1,\quad A_1+a_1\le1.
\]
Take $2n$ mutually independent variables $U_{0i},U_{1i}$ with these laws,
and define
\[
 T_j=j+\sum_{i\le j}U_{1i}+\sum_{i>j}U_{0i},\quad 0\le j\le n.
\]
If $na_0\le t<n(1-a_1)$, there is a unique $K\in\{0,\ldots,n-1\}$
with $T_K\le t<T_{K+1}$. Put
\[
 b=\min\{1-A_1-a_1,1-a_0-a_1\}\ge0,\qquad R=t-T_K.
\]
Then
\begin{equation}\label{eq:gain}
              g_n(t):=t-\E K\ \ge\ \E\min\{R,b\}\ \ge0.
\end{equation}
More precisely, let $D=1+U_1$, $F_0(x)=\Pp(U_0\le x)$,
$M_0(x)=\E[U_0\mathbf1_{\{U_0\le x\}}]$, and
\begin{align*}
 h(x)&=F_0(x)\{\Pp(D>x)-\E(D-x)_+\}-M_0(x),\\
 R_i&=t-(i-1)-\sum_{l<i}U_{1l}-\sum_{l>i}U_{0l}.
\end{align*}
The exact cancellation identity is
\begin{equation}\label{eq:identity}
                       g_n(t)=\sum_{i=1}^n\E h(R_i).
\end{equation}
\end{theorem}
\begin{proof}
Each increment $T_i-T_{i-1}=1+U_{1i}-U_{0i}$ is at least
$1-a_0-a_1>0$. The threshold assumptions give $T_0\le t<T_n$.
Thus the closed lower and strict upper crossing events partition the
sample space, even when $t$ is an atom of a conditional sum.

Write $W_j=T_j-j$ and $O=T_{K+1}-t$. The identity
$T_{K+1}=K+1+W_{K+1}$ gives
\[
                   t-\E K=1+\E W_{K+1}-\E O.               \tag{3}
\]
Expand the selected noise sum as
\[
 W_{K+1}=\sum_{i=1}^n U_{0i}
       +\sum_{i=1}^n(U_{1i}-U_{0i})\mathbf1_{\{T_{i-1}\le t\}}.
\]
The event in the $i$th summand is $\{U_{0i}\le R_i\}$.
Both $R_i$ and $U_{0i}$ are independent of the centered $U_{1i}$, so
\[
                      \E W_{K+1}=-\sum_i\E M_0(R_i).       \tag{4}
\]
It is generally positive, not zero. A naive optional-stopping argument
would lose this correction.

Condition on $R_i=x$. Crossing at $i$ means $U_{0i}\le x<D_i$.
Independence of the pair gives crossing probability
$F_0(x)\Pp(D>x)$ and the overshoot contribution
$F_0(x)\E(D-x)_+$. Sum these quantities in (3), using that crossing
probabilities sum to one, and use (4). This proves \eqref{eq:identity}.

Set $m=1-a_1$. If $x<m$, then $D>x$ almost surely and $\E D=1$;
therefore
\[
                         h(x)=\E(x-U_0)_+\ge0.             \tag{5}
\]
In this case it is also at least the conditional expected value of
$\min\{x-U_0,b\}$ on the crossing event. If $x\ge m$, then
$x>a_0$, hence $F_0(x)=1$ and $M_0(x)=0$. Since $D\le1+A_1$,
\[
 h(x)=\Pp(D>x)-\E(D-x)_+
       \ge(1-A_1-a_1)\Pp(D>x)\ge b\Pp(D>x).               \tag{6}
\]
At a crossing in this second case the age is $x-U_0\ge m-a_0\ge b$.
Consequently (5)--(6), summed over the unique crossing, give
\eqref{eq:gain}. Formula (6) is used at $x=m$, so atoms at the lower
endpoint of $D$ introduce no change of convention. All selected sums are
integrable, and every sum in the proof is finite.
\end{proof}

\section*{An exact replacement for the separate count sums}
For $Y=B+U$ with $\Pp(B=1)=p$ and these conditional laws, the conditional
distribution of a raw iid sum given $j$ successes is precisely that of
$T_j$. With $F_B$ the $\mathrm{Bin}(n,p)$ CDF,
\begin{equation}\label{eq:cdf}
            \Pp(Y_1+\cdots+Y_n\le t)=\E F_B(K).
\end{equation}
The coupling is auxiliary: it preserves each conditional law, hence
preserves the actual iid mixture. It does not impose independence between
$K$ and the binomial count or substitute an optimized relaxation.

For the acquired right strip, $a_0+a_1\le.072$ and $A_1=.6$, one may
take $b=.328$ in \eqref{eq:gain}; for the left strip, inward sum at most
$.068$ and $A_1=.7$ give $b=.232$. For $n\ge42$, $p\in[.25,.53]$,
$s=\E U^2$, $ns\le.616$ and $|H|\le1.84$, put
$t=np+H\sqrt{p(1-p)+s}$. Then
\[
 .25n-.948<t<.53n+.948,\qquad
                   na_0\le t<n(1-a_1).
\]
Indeed $\sqrt{1/4+.616/42}<.515$ and $1.84(.515)<.948$;
the remaining two inequalities hold already at $n=42$. Thus the
cancellation identity covers the whole retained acquired noise range.
It makes no assertion that the acquisition hypotheses themselves hold
without their existing contact and predecessor dependencies.

To state exactly what remains, put $v=p(1-p)$,
$\rho=\E|Y-p|^3/(v+s)^{3/2}$, $\rho_B=(1-2v)/\sqrt v$, and $t=k+u$.
For $\pi_k=\Pp(B_n=k)$ and
$Q_k(j)=F_B(j)-F_B(k)-\pi_k(j-k)$, the direct comparison is exactly
\begin{align}\label{eq:delta}
 \Delta_C={}&-\sqrt n\,\pi_k g_n(t)+\sqrt n\,\E Q_k(K)
              +\sqrt n\,\pi_k u\notag\\
 &-\sqrt n\left[
 \Phi\!\left(\frac{t-np}{\sqrt{n(v+s)}}\right)
 -\Phi\!\left(\frac{k-np}{\sqrt{nv}}\right)\right]
              +C(\rho_B-\rho).
\end{align}
The leading count contribution in this formula has the correct sign, for
all phases, with both conditional noise laws present. This is stronger than
an asymptotic assertion: every identity is exact for the actual sample size.
But neither $\E Q_k(K)$ nor the two repairs in \eqref{eq:delta} has been
absorbed by \eqref{eq:gain}. That joint inequality is the remaining task.
The current general count certificates therefore remain dependencies of the
accepted comparison. They cannot be removed merely because the leading
term has been simplified.

\section*{Controls and limitations}
The companion exact checks reconstruct all conditional sums and the local
kernel $h$ separately, including closed collisions and strict crossings.
They check \eqref{eq:identity}, \eqref{eq:gain}, \eqref{eq:cdf} and the
nonzero selected-noise bias. They also include a false extension: take
$U_0=0$ and $D=9/10,39/20$ with probabilities $19/21,2/21$.
At $n=2$, $t=9/10$, $\E K-t=1/210>0$; here $A_1+a_1=21/20>1$.
No assertion of strict gain proportional to $s$, or of the full sharp
Berry--Esseen bound, follows from this note alone.
\end{document}
